\QHeading{第1题　填一填。}

\par\noindent\makebox[\linewidth][c]{
\begin{tikzpicture}[x=1bp,y=-1bp,color=black,draw=black,text=black]
\begin{scope}\clip (56.0,113.0) rectangle (401.0,244.0);
\hdhead{45}{多边形的认识——信息窗3}
\hdsection{1}{132}{填一填。}
\begin{scope}[scale=.6666667,draw=black,text=black]\draw[hdoutline] (382.000,308.000)--(419.000,203.000)--(513.000,203.000)--(550.000,308.000)--cycle;\end{scope}
\begin{scope}[scale=.6666667,draw=black,text=black]\draw[hdoutline,dashed] (419.000,203.000)--(419.000,308.000);\end{scope}
\begin{scope}[scale=.6666667,draw=black,text=black]\draw[line width=.45pt] (419.000,298.000)--(429.000,298.000)--(429.000,308.000);\end{scope}
\q{423}{178}{28}{13}{（}
\q{480}{178}{25}{13}{）}
\q{332}{249}{28}{13}{（}
\q{384}{249}{25}{13}{）}
\q{418}{249}{28}{13}{（}
\q{465}{249}{25}{13}{）}
\q{529}{249}{28}{13}{（}
\q{587}{249}{25}{13}{）}
\q{420}{315}{28}{13}{（}
\q{482}{315}{25}{13}{）}
\hdsection{2}{369}{分别画出下面梯形底边上的高。}
\begin{scope}[scale=.6666667,draw=black,text=black]\draw[hdoutline] (124.000,498.000)--(168.000,428.000)--(261.000,428.000)--(275.000,498.000)--cycle;\end{scope}
\begin{scope}[scale=.6666667,draw=black,text=black]\draw[hdoutline] (356.000,413.000)--(424.000,458.000)--(424.000,520.000)--(356.000,520.000)--cycle;\end{scope}
\begin{scope}[scale=.6666667,draw=black,text=black]\draw[hdoutline] (507.000,447.000)--(634.000,447.000)--(608.000,517.000)--(534.000,517.000)--cycle;\end{scope}
\begin{scope}[scale=.6666667,draw=black,text=black]\draw[hdoutline] (734.000,440.000)--(803.000,440.000)--(803.000,517.000)--(734.000,492.000)--cycle;\end{scope}
\begin{scope}[scale=.6666667,draw=black,text=black]\node[inner sep=0pt,font=\normalsize,text=black] at (201.000,516.000) {底};\end{scope}
\begin{scope}[scale=.6666667,draw=black,text=black]\node[inner sep=0pt,font=\normalsize,text=black] at (344.000,471.000) {底};\end{scope}
\begin{scope}[scale=.6666667,draw=black,text=black]\node[inner sep=0pt,font=\normalsize,text=black] at (571.000,431.000) {底};\end{scope}
\begin{scope}[scale=.6666667,draw=black,text=black]\node[inner sep=0pt,font=\normalsize,text=black] at (722.000,468.000) {底};\end{scope}
\hdsection{3}{559}{画一画，分一分。}
\q{102}{607}{730}{12}{① 在图形中各画一条线段，把图形分成一个平行四边形和一个三角形。}
\begin{scope}[scale=.6666667,draw=black,text=black]\draw[hdoutline] (300.000,715.000)--(322.000,647.000)--(399.000,647.000)--(440.000,715.000)--cycle;\end{scope}
\begin{scope}[scale=.6666667,draw=black,text=black]\draw[hdoutline] (559.000,648.000)--(685.000,648.000)--(642.000,715.000)--(583.000,715.000)--cycle;\end{scope}
\q{103}{738}{734}{12}{② 在每个图形中各画一条线段，把图形分成两个完全相同的梯形。}
\begin{scope}[scale=.6666667,draw=black,text=black]\draw[hdoutline] (301.500,846.000)--(339.500,776.000)--(447.500,776.000)--(409.500,846.000)--cycle;\end{scope}
\begin{scope}[scale=.6666667,draw=black,text=black]\draw[hdoutline] (571.000,776.000)--(680.000,776.000)--(680.000,846.000)--(571.000,846.000)--cycle;\end{scope}
\hdsection{4}{879}{在下面的点子图中分别画一个长方形、一个正方形、一个平行四边形和一个等腰梯形。}
\begin{scope}[scale=.6666667,draw=black,text=black]\foreach \i in {0,...,38}{\foreach \j in {0,...,9}{\fill[black] ({113.6+\i*18.4},{974.2+\j*18.4}) circle (1.15);}}\end{scope}

\begin{scope}[color=answerred,draw=answerred,text=answerred]\ans{296.67}{118.67}{46.67}{10}{上底}
\ans{294.67}{210.00}{46.67}{10}{下底}
\ans{236.00}{166.00}{46.67}{10}{腰}
\ans{367.33}{166.00}{46.67}{10}{腰}
\ans{293.33}{166.00}{46.67}{10}{高}
\draw[answerred,line width=.7pt,dashed] (112.000,285.333) -- (112.000,332.000);
\draw[answerred,line width=.7pt] (117.333,332.000) -- (117.333,326.667);
\draw[answerred,line width=.7pt] (117.333,326.667) -- (112.000,326.667);
\ans{116.00}{303.33}{20.00}{9}{高}
\draw[answerred,line width=.7pt,dashed] (282.667,305.333) -- (237.333,305.333);
\draw[answerred,line width=.7pt] (237.333,310.667) -- (242.667,310.667);
\draw[answerred,line width=.7pt] (242.667,310.667) -- (242.667,305.333);
\ans{264.00}{300.00}{20.00}{9}{高}
\draw[answerred,line width=.7pt,dashed] (356.000,344.667) -- (356.000,298.000);
\draw[answerred,line width=.7pt] (361.333,298.000) -- (361.333,303.333);
\draw[answerred,line width=.7pt] (361.333,303.333) -- (356.000,303.333);
\ans{360.00}{316.00}{20.00}{9}{高}
\draw[answerred,line width=.7pt,dashed] (535.333,306.667) -- (489.333,306.667);
\draw[answerred,line width=.7pt] (489.333,312.000) -- (494.667,312.000);
\draw[answerred,line width=.7pt] (494.667,312.000) -- (494.667,306.667);
\ans{506.67}{309.33}{20.00}{9}{高}
\draw[answerred,line width=.7pt] (266.000,431.333) -- (251.333,476.667);
\draw[answerred,line width=.7pt] (428.000,476.667) -- (412.000,432.000);
\draw[answerred,line width=.7pt] (243.333,517.333) -- (256.000,564.000);
\draw[answerred,line width=.7pt] (400.000,517.333) -- (433.333,564.000);
\draw[answerred,line width=.8pt] (88.000,735.333) -- (174.000,735.333) -- (174.000,674.000) -- (88.000,674.000) -- cycle;
\draw[answerred,line width=.8pt] (211.333,735.333) -- (272.667,735.333) -- (272.667,674.000) -- (211.333,674.000) -- cycle;
\draw[answerred,line width=.8pt] (309.333,735.333) -- (371.333,735.333) -- (396.000,674.000) -- (334.000,674.000) -- cycle;
\draw[answerred,line width=.8pt] (433.333,735.333) -- (532.667,735.333) -- (507.333,686.000) -- (458.667,686.000) -- cycle;
\ans{100.00}{750.67}{100.00}{9}{长方形}
\ans{220.67}{750.67}{100.00}{9}{正方形}
\ans{319.33}{750.67}{100.00}{9}{平行四边形}
\ans{453.33}{750.67}{100.00}{9}{等腰梯形}
\end{scope}
\end{scope}\pgfresetboundingbox\path[use as bounding box] (56.0,113.0) rectangle (401.0,244.0);
\end{tikzpicture}}\par
\par\noindent




\par

\par\noindent

\QHeading{第2题　分别画出下面梯形底边上的高。}

\par\noindent\makebox[\linewidth][c]{
\begin{tikzpicture}[x=1bp,y=-1bp,color=black,draw=black,text=black]
\begin{scope}\clip (56.0,270.0) rectangle (541.0,370.67);
\hdhead{45}{多边形的认识——信息窗3}
\hdsection{1}{132}{填一填。}
\begin{scope}[scale=.6666667,draw=black,text=black]\draw[hdoutline] (382.000,308.000)--(419.000,203.000)--(513.000,203.000)--(550.000,308.000)--cycle;\end{scope}
\begin{scope}[scale=.6666667,draw=black,text=black]\draw[hdoutline,dashed] (419.000,203.000)--(419.000,308.000);\end{scope}
\begin{scope}[scale=.6666667,draw=black,text=black]\draw[line width=.45pt] (419.000,298.000)--(429.000,298.000)--(429.000,308.000);\end{scope}
\q{423}{178}{28}{13}{（}
\q{480}{178}{25}{13}{）}
\q{332}{249}{28}{13}{（}
\q{384}{249}{25}{13}{）}
\q{418}{249}{28}{13}{（}
\q{465}{249}{25}{13}{）}
\q{529}{249}{28}{13}{（}
\q{587}{249}{25}{13}{）}
\q{420}{315}{28}{13}{（}
\q{482}{315}{25}{13}{）}
\hdsection{2}{369}{分别画出下面梯形底边上的高。}
\begin{scope}[scale=.6666667,draw=black,text=black]\draw[hdoutline] (124.000,498.000)--(168.000,428.000)--(261.000,428.000)--(275.000,498.000)--cycle;\end{scope}
\begin{scope}[scale=.6666667,draw=black,text=black]\draw[hdoutline] (356.000,413.000)--(424.000,458.000)--(424.000,520.000)--(356.000,520.000)--cycle;\end{scope}
\begin{scope}[scale=.6666667,draw=black,text=black]\draw[hdoutline] (507.000,447.000)--(634.000,447.000)--(608.000,517.000)--(534.000,517.000)--cycle;\end{scope}
\begin{scope}[scale=.6666667,draw=black,text=black]\draw[hdoutline] (734.000,440.000)--(803.000,440.000)--(803.000,517.000)--(734.000,492.000)--cycle;\end{scope}
\begin{scope}[scale=.6666667,draw=black,text=black]\node[inner sep=0pt,font=\normalsize,text=black] at (201.000,516.000) {底};\end{scope}
\begin{scope}[scale=.6666667,draw=black,text=black]\node[inner sep=0pt,font=\normalsize,text=black] at (344.000,471.000) {底};\end{scope}
\begin{scope}[scale=.6666667,draw=black,text=black]\node[inner sep=0pt,font=\normalsize,text=black] at (571.000,431.000) {底};\end{scope}
\begin{scope}[scale=.6666667,draw=black,text=black]\node[inner sep=0pt,font=\normalsize,text=black] at (722.000,468.000) {底};\end{scope}
\hdsection{3}{559}{画一画，分一分。}
\q{102}{607}{730}{12}{① 在图形中各画一条线段，把图形分成一个平行四边形和一个三角形。}
\begin{scope}[scale=.6666667,draw=black,text=black]\draw[hdoutline] (300.000,715.000)--(322.000,647.000)--(399.000,647.000)--(440.000,715.000)--cycle;\end{scope}
\begin{scope}[scale=.6666667,draw=black,text=black]\draw[hdoutline] (559.000,648.000)--(685.000,648.000)--(642.000,715.000)--(583.000,715.000)--cycle;\end{scope}
\q{103}{738}{734}{12}{② 在每个图形中各画一条线段，把图形分成两个完全相同的梯形。}
\begin{scope}[scale=.6666667,draw=black,text=black]\draw[hdoutline] (301.500,846.000)--(339.500,776.000)--(447.500,776.000)--(409.500,846.000)--cycle;\end{scope}
\begin{scope}[scale=.6666667,draw=black,text=black]\draw[hdoutline] (571.000,776.000)--(680.000,776.000)--(680.000,846.000)--(571.000,846.000)--cycle;\end{scope}
\hdsection{4}{879}{在下面的点子图中分别画一个长方形、一个正方形、一个平行四边形和一个等腰梯形。}
\begin{scope}[scale=.6666667,draw=black,text=black]\foreach \i in {0,...,38}{\foreach \j in {0,...,9}{\fill[black] ({113.6+\i*18.4},{974.2+\j*18.4}) circle (1.15);}}\end{scope}

\begin{scope}[color=answerred,draw=answerred,text=answerred]\ans{296.67}{118.67}{46.67}{10}{上底}
\ans{294.67}{210.00}{46.67}{10}{下底}
\ans{236.00}{166.00}{46.67}{10}{腰}
\ans{367.33}{166.00}{46.67}{10}{腰}
\ans{293.33}{166.00}{46.67}{10}{高}
\draw[answerred,line width=.7pt,dashed] (112.000,285.333) -- (112.000,332.000);
\draw[answerred,line width=.7pt] (117.333,332.000) -- (117.333,326.667);
\draw[answerred,line width=.7pt] (117.333,326.667) -- (112.000,326.667);
\ans{116.00}{303.33}{20.00}{9}{高}
\draw[answerred,line width=.7pt,dashed] (282.667,305.333) -- (237.333,305.333);
\draw[answerred,line width=.7pt] (237.333,310.667) -- (242.667,310.667);
\draw[answerred,line width=.7pt] (242.667,310.667) -- (242.667,305.333);
\ans{264.00}{300.00}{20.00}{9}{高}
\draw[answerred,line width=.7pt,dashed] (356.000,344.667) -- (356.000,298.000);
\draw[answerred,line width=.7pt] (361.333,298.000) -- (361.333,303.333);
\draw[answerred,line width=.7pt] (361.333,303.333) -- (356.000,303.333);
\ans{360.00}{316.00}{20.00}{9}{高}
\draw[answerred,line width=.7pt,dashed] (535.333,306.667) -- (489.333,306.667);
\draw[answerred,line width=.7pt] (489.333,312.000) -- (494.667,312.000);
\draw[answerred,line width=.7pt] (494.667,312.000) -- (494.667,306.667);
\ans{506.67}{309.33}{20.00}{9}{高}
\draw[answerred,line width=.7pt] (266.000,431.333) -- (251.333,476.667);
\draw[answerred,line width=.7pt] (428.000,476.667) -- (412.000,432.000);
\draw[answerred,line width=.7pt] (243.333,517.333) -- (256.000,564.000);
\draw[answerred,line width=.7pt] (400.000,517.333) -- (433.333,564.000);
\draw[answerred,line width=.8pt] (88.000,735.333) -- (174.000,735.333) -- (174.000,674.000) -- (88.000,674.000) -- cycle;
\draw[answerred,line width=.8pt] (211.333,735.333) -- (272.667,735.333) -- (272.667,674.000) -- (211.333,674.000) -- cycle;
\draw[answerred,line width=.8pt] (309.333,735.333) -- (371.333,735.333) -- (396.000,674.000) -- (334.000,674.000) -- cycle;
\draw[answerred,line width=.8pt] (433.333,735.333) -- (532.667,735.333) -- (507.333,686.000) -- (458.667,686.000) -- cycle;
\ans{100.00}{750.67}{100.00}{9}{长方形}
\ans{220.67}{750.67}{100.00}{9}{正方形}
\ans{319.33}{750.67}{100.00}{9}{平行四边形}
\ans{453.33}{750.67}{100.00}{9}{等腰梯形}
\end{scope}
\end{scope}\pgfresetboundingbox\path[use as bounding box] (56.0,270.0) rectangle (541.0,370.67);
\end{tikzpicture}}\par



\par







\par\noindent

\QHeading{第3题　画一画，分一分。}

\par\noindent\makebox[\linewidth][c]{
\begin{tikzpicture}[x=1bp,y=-1bp,color=black,draw=black,text=black]
\begin{scope}\clip (56.0,512.0) rectangle (459.0,584.0);
\hdhead{45}{多边形的认识——信息窗3}
\hdsection{1}{132}{填一填。}
\begin{scope}[scale=.6666667,draw=black,text=black]\draw[hdoutline] (382.000,308.000)--(419.000,203.000)--(513.000,203.000)--(550.000,308.000)--cycle;\end{scope}
\begin{scope}[scale=.6666667,draw=black,text=black]\draw[hdoutline,dashed] (419.000,203.000)--(419.000,308.000);\end{scope}
\begin{scope}[scale=.6666667,draw=black,text=black]\draw[line width=.45pt] (419.000,298.000)--(429.000,298.000)--(429.000,308.000);\end{scope}
\q{423}{178}{28}{13}{（}
\q{480}{178}{25}{13}{）}
\q{332}{249}{28}{13}{（}
\q{384}{249}{25}{13}{）}
\q{418}{249}{28}{13}{（}
\q{465}{249}{25}{13}{）}
\q{529}{249}{28}{13}{（}
\q{587}{249}{25}{13}{）}
\q{420}{315}{28}{13}{（}
\q{482}{315}{25}{13}{）}
\hdsection{2}{369}{分别画出下面梯形底边上的高。}
\begin{scope}[scale=.6666667,draw=black,text=black]\draw[hdoutline] (124.000,498.000)--(168.000,428.000)--(261.000,428.000)--(275.000,498.000)--cycle;\end{scope}
\begin{scope}[scale=.6666667,draw=black,text=black]\draw[hdoutline] (356.000,413.000)--(424.000,458.000)--(424.000,520.000)--(356.000,520.000)--cycle;\end{scope}
\begin{scope}[scale=.6666667,draw=black,text=black]\draw[hdoutline] (507.000,447.000)--(634.000,447.000)--(608.000,517.000)--(534.000,517.000)--cycle;\end{scope}
\begin{scope}[scale=.6666667,draw=black,text=black]\draw[hdoutline] (734.000,440.000)--(803.000,440.000)--(803.000,517.000)--(734.000,492.000)--cycle;\end{scope}
\begin{scope}[scale=.6666667,draw=black,text=black]\node[inner sep=0pt,font=\normalsize,text=black] at (201.000,516.000) {底};\end{scope}
\begin{scope}[scale=.6666667,draw=black,text=black]\node[inner sep=0pt,font=\normalsize,text=black] at (344.000,471.000) {底};\end{scope}
\begin{scope}[scale=.6666667,draw=black,text=black]\node[inner sep=0pt,font=\normalsize,text=black] at (571.000,431.000) {底};\end{scope}
\begin{scope}[scale=.6666667,draw=black,text=black]\node[inner sep=0pt,font=\normalsize,text=black] at (722.000,468.000) {底};\end{scope}
\hdsection{3}{559}{画一画，分一分。}

\begin{scope}[scale=.6666667,draw=black,text=black]\draw[hdoutline] (300.000,715.000)--(322.000,647.000)--(399.000,647.000)--(440.000,715.000)--cycle;\end{scope}
\begin{scope}[scale=.6666667,draw=black,text=black]\draw[hdoutline] (559.000,648.000)--(685.000,648.000)--(642.000,715.000)--(583.000,715.000)--cycle;\end{scope}

\begin{scope}[scale=.6666667,draw=black,text=black]\draw[hdoutline] (301.500,846.000)--(339.500,776.000)--(447.500,776.000)--(409.500,846.000)--cycle;\end{scope}
\begin{scope}[scale=.6666667,draw=black,text=black]\draw[hdoutline] (571.000,776.000)--(680.000,776.000)--(680.000,846.000)--(571.000,846.000)--cycle;\end{scope}
\hdsection{4}{879}{在下面的点子图中分别画一个长方形、一个正方形、一个平行四边形和一个等腰梯形。}
\begin{scope}[scale=.6666667,draw=black,text=black]\foreach \i in {0,...,38}{\foreach \j in {0,...,9}{\fill[black] ({113.6+\i*18.4},{974.2+\j*18.4}) circle (1.15);}}\end{scope}

\begin{scope}[color=answerred,draw=answerred,text=answerred]\ans{296.67}{118.67}{46.67}{10}{上底}
\ans{294.67}{210.00}{46.67}{10}{下底}
\ans{236.00}{166.00}{46.67}{10}{腰}
\ans{367.33}{166.00}{46.67}{10}{腰}
\ans{293.33}{166.00}{46.67}{10}{高}
\draw[answerred,line width=.7pt,dashed] (112.000,285.333) -- (112.000,332.000);
\draw[answerred,line width=.7pt] (117.333,332.000) -- (117.333,326.667);
\draw[answerred,line width=.7pt] (117.333,326.667) -- (112.000,326.667);
\ans{116.00}{303.33}{20.00}{9}{高}
\draw[answerred,line width=.7pt,dashed] (282.667,305.333) -- (237.333,305.333);
\draw[answerred,line width=.7pt] (237.333,310.667) -- (242.667,310.667);
\draw[answerred,line width=.7pt] (242.667,310.667) -- (242.667,305.333);
\ans{264.00}{300.00}{20.00}{9}{高}
\draw[answerred,line width=.7pt,dashed] (356.000,344.667) -- (356.000,298.000);
\draw[answerred,line width=.7pt] (361.333,298.000) -- (361.333,303.333);
\draw[answerred,line width=.7pt] (361.333,303.333) -- (356.000,303.333);
\ans{360.00}{316.00}{20.00}{9}{高}
\draw[answerred,line width=.7pt,dashed] (535.333,306.667) -- (489.333,306.667);
\draw[answerred,line width=.7pt] (489.333,312.000) -- (494.667,312.000);
\draw[answerred,line width=.7pt] (494.667,312.000) -- (494.667,306.667);
\ans{506.67}{309.33}{20.00}{9}{高}
\draw[answerred,line width=.7pt] (266.000,431.333) -- (251.333,476.667);
\draw[answerred,line width=.7pt] (428.000,476.667) -- (412.000,432.000);
\draw[answerred,line width=.7pt] (243.333,517.333) -- (256.000,564.000);
\draw[answerred,line width=.7pt] (400.000,517.333) -- (433.333,564.000);
\draw[answerred,line width=.8pt] (88.000,735.333) -- (174.000,735.333) -- (174.000,674.000) -- (88.000,674.000) -- cycle;
\draw[answerred,line width=.8pt] (211.333,735.333) -- (272.667,735.333) -- (272.667,674.000) -- (211.333,674.000) -- cycle;
\draw[answerred,line width=.8pt] (309.333,735.333) -- (371.333,735.333) -- (396.000,674.000) -- (334.000,674.000) -- cycle;
\draw[answerred,line width=.8pt] (433.333,735.333) -- (532.667,735.333) -- (507.333,686.000) -- (458.667,686.000) -- cycle;
\ans{100.00}{750.67}{100.00}{9}{长方形}
\ans{220.67}{750.67}{100.00}{9}{正方形}
\ans{319.33}{750.67}{100.00}{9}{平行四边形}
\ans{453.33}{750.67}{100.00}{9}{等腰梯形}
\end{scope}
\end{scope}\pgfresetboundingbox\path[use as bounding box] (56.0,512.0) rectangle (459.0,584.0);
\end{tikzpicture}}\par

\PlainItem{① 在图形中各画一条线段，把图形分成一个平行四边形和一个三角形。}
\par



\PlainItem{② 在每个图形中各画一条线段，把图形分成两个完全相同的梯形。}


\par\noindent

\QHeading{第3题①　从另一端分割梯形}

\Answer{在第一幅梯形中，除过右上顶点作平行于左腰的分割线外，也可以过左上顶点作平行于右腰的分割线。倒放的梯形同理，可从另一端作。}
\par\smallskip\noindent\begin{adjustbox}{max width=\linewidth,center}
\begin{tikzpicture}[x=.035cm,y=-.035cm,font=\normalsize]
\begin{scope}[shift={(-300,-647)}]
\draw[black,thick] (300,715)--(322,647)--(399,647)--(440,715)--cycle;
\draw[answerred,thick] (322,647)--(363,715);
\end{scope}
\begin{scope}[xshift=5.5cm,shift={(-559,-648)}]
\draw[black,thick] (559,648)--(685,648)--(642,715)--(583,715)--cycle;
\draw[answerred,thick] (583,715)--(626,648);
\end{scope}
\end{tikzpicture}
\end{adjustbox}\par\smallskip
\Answer{新增红色线段各把原梯形分成一个平行四边形和一个三角形。分割线用实线；它们不是作高线。}

\QHeading{第3题②　分成两个全等梯形的另一些画法}

\Answer{经过图形中心，连接一对关于中心相对称的上下底上的点，可以得到两块完全相同的图形；为得到梯形，分割线不能与原来的腰平行。}
\par\smallskip\noindent\begin{adjustbox}{max width=\linewidth,center}
\begin{tikzpicture}[x=.65cm,y=.65cm,font=\normalsize]
\draw[black,thick] (0,0)--(4,0)--(5,2)--(1,2)--cycle;
\draw[answerred,thick] (1.2,0)--(3.8,2);
\fill[answerred] (2.5,1) circle (1.4pt);
\begin{scope}[xshift=4.4cm]
\draw[black,thick] (0,0) rectangle (4,2);
\draw[answerred,thick] (1,0)--(3,2);
\fill[answerred] (2,1) circle (1.4pt);
\end{scope}
\end{tikzpicture}
\end{adjustbox}\par\smallskip
\Answer{两部分可绕中心旋转180度后重合。矩形不能简单从中间画一条竖线来代替本题答案，否则分出的两块都是长方形，不是本题“只有一组对边平行”的梯形。}

\QHeading{第4题　在下面的点子图中分别画一个长方形、一个正方形、一个平行四边形和一个等腰梯形。}

\par\noindent\makebox[\linewidth][c]{
\begin{tikzpicture}[x=1bp,y=-1bp,color=black,draw=black,text=black]
\begin{scope}\clip (70.0,608.0) rectangle (548.0,766.0);
\hdhead{45}{多边形的认识——信息窗3}
\hdsection{1}{132}{填一填。}
\begin{scope}[scale=.6666667,draw=black,text=black]\draw[hdoutline] (382.000,308.000)--(419.000,203.000)--(513.000,203.000)--(550.000,308.000)--cycle;\end{scope}
\begin{scope}[scale=.6666667,draw=black,text=black]\draw[hdoutline,dashed] (419.000,203.000)--(419.000,308.000);\end{scope}
\begin{scope}[scale=.6666667,draw=black,text=black]\draw[line width=.45pt] (419.000,298.000)--(429.000,298.000)--(429.000,308.000);\end{scope}
\q{423}{178}{28}{13}{（}
\q{480}{178}{25}{13}{）}
\q{332}{249}{28}{13}{（}
\q{384}{249}{25}{13}{）}
\q{418}{249}{28}{13}{（}
\q{465}{249}{25}{13}{）}
\q{529}{249}{28}{13}{（}
\q{587}{249}{25}{13}{）}
\q{420}{315}{28}{13}{（}
\q{482}{315}{25}{13}{）}
\hdsection{2}{369}{分别画出下面梯形底边上的高。}
\begin{scope}[scale=.6666667,draw=black,text=black]\draw[hdoutline] (124.000,498.000)--(168.000,428.000)--(261.000,428.000)--(275.000,498.000)--cycle;\end{scope}
\begin{scope}[scale=.6666667,draw=black,text=black]\draw[hdoutline] (356.000,413.000)--(424.000,458.000)--(424.000,520.000)--(356.000,520.000)--cycle;\end{scope}
\begin{scope}[scale=.6666667,draw=black,text=black]\draw[hdoutline] (507.000,447.000)--(634.000,447.000)--(608.000,517.000)--(534.000,517.000)--cycle;\end{scope}
\begin{scope}[scale=.6666667,draw=black,text=black]\draw[hdoutline] (734.000,440.000)--(803.000,440.000)--(803.000,517.000)--(734.000,492.000)--cycle;\end{scope}
\begin{scope}[scale=.6666667,draw=black,text=black]\node[inner sep=0pt,font=\normalsize,text=black] at (201.000,516.000) {底};\end{scope}
\begin{scope}[scale=.6666667,draw=black,text=black]\node[inner sep=0pt,font=\normalsize,text=black] at (344.000,471.000) {底};\end{scope}
\begin{scope}[scale=.6666667,draw=black,text=black]\node[inner sep=0pt,font=\normalsize,text=black] at (571.000,431.000) {底};\end{scope}
\begin{scope}[scale=.6666667,draw=black,text=black]\node[inner sep=0pt,font=\normalsize,text=black] at (722.000,468.000) {底};\end{scope}
\hdsection{3}{559}{画一画，分一分。}
\q{102}{607}{730}{12}{① 在图形中各画一条线段，把图形分成一个平行四边形和一个三角形。}
\begin{scope}[scale=.6666667,draw=black,text=black]\draw[hdoutline] (300.000,715.000)--(322.000,647.000)--(399.000,647.000)--(440.000,715.000)--cycle;\end{scope}
\begin{scope}[scale=.6666667,draw=black,text=black]\draw[hdoutline] (559.000,648.000)--(685.000,648.000)--(642.000,715.000)--(583.000,715.000)--cycle;\end{scope}
\q{103}{738}{734}{12}{② 在每个图形中各画一条线段，把图形分成两个完全相同的梯形。}
\begin{scope}[scale=.6666667,draw=black,text=black]\draw[hdoutline] (301.500,846.000)--(339.500,776.000)--(447.500,776.000)--(409.500,846.000)--cycle;\end{scope}
\begin{scope}[scale=.6666667,draw=black,text=black]\draw[hdoutline] (571.000,776.000)--(680.000,776.000)--(680.000,846.000)--(571.000,846.000)--cycle;\end{scope}
\hdsection{4}{879}{在下面的点子图中分别画一个长方形、一个正方形、一个平行四边形和一个等腰梯形。}
\begin{scope}[scale=.6666667,draw=black,text=black]\foreach \i in {0,...,38}{\foreach \j in {0,...,9}{\fill[black] ({113.6+\i*18.4},{974.2+\j*18.4}) circle (1.15);}}\end{scope}

\begin{scope}[color=answerred,draw=answerred,text=answerred]\ans{296.67}{118.67}{46.67}{10}{上底}
\ans{294.67}{210.00}{46.67}{10}{下底}
\ans{236.00}{166.00}{46.67}{10}{腰}
\ans{367.33}{166.00}{46.67}{10}{腰}
\ans{293.33}{166.00}{46.67}{10}{高}
\draw[answerred,line width=.7pt,dashed] (112.000,285.333) -- (112.000,332.000);
\draw[answerred,line width=.7pt] (117.333,332.000) -- (117.333,326.667);
\draw[answerred,line width=.7pt] (117.333,326.667) -- (112.000,326.667);
\ans{116.00}{303.33}{20.00}{9}{高}
\draw[answerred,line width=.7pt,dashed] (282.667,305.333) -- (237.333,305.333);
\draw[answerred,line width=.7pt] (237.333,310.667) -- (242.667,310.667);
\draw[answerred,line width=.7pt] (242.667,310.667) -- (242.667,305.333);
\ans{264.00}{300.00}{20.00}{9}{高}
\draw[answerred,line width=.7pt,dashed] (356.000,344.667) -- (356.000,298.000);
\draw[answerred,line width=.7pt] (361.333,298.000) -- (361.333,303.333);
\draw[answerred,line width=.7pt] (361.333,303.333) -- (356.000,303.333);
\ans{360.00}{316.00}{20.00}{9}{高}
\draw[answerred,line width=.7pt,dashed] (535.333,306.667) -- (489.333,306.667);
\draw[answerred,line width=.7pt] (489.333,312.000) -- (494.667,312.000);
\draw[answerred,line width=.7pt] (494.667,312.000) -- (494.667,306.667);
\ans{506.67}{309.33}{20.00}{9}{高}
\draw[answerred,line width=.7pt] (266.000,431.333) -- (251.333,476.667);
\draw[answerred,line width=.7pt] (428.000,476.667) -- (412.000,432.000);
\draw[answerred,line width=.7pt] (243.333,517.333) -- (256.000,564.000);
\draw[answerred,line width=.7pt] (400.000,517.333) -- (433.333,564.000);
\draw[answerred,line width=.8pt] (88.000,735.333) -- (174.000,735.333) -- (174.000,674.000) -- (88.000,674.000) -- cycle;
\draw[answerred,line width=.8pt] (211.333,735.333) -- (272.667,735.333) -- (272.667,674.000) -- (211.333,674.000) -- cycle;
\draw[answerred,line width=.8pt] (309.333,735.333) -- (371.333,735.333) -- (396.000,674.000) -- (334.000,674.000) -- cycle;
\draw[answerred,line width=.8pt] (433.333,735.333) -- (532.667,735.333) -- (507.333,686.000) -- (458.667,686.000) -- cycle;
\ans{100.00}{750.67}{100.00}{9}{长方形}
\ans{220.67}{750.67}{100.00}{9}{正方形}
\ans{319.33}{750.67}{100.00}{9}{平行四边形}
\ans{453.33}{750.67}{100.00}{9}{等腰梯形}
\end{scope}
\end{scope}\pgfresetboundingbox\path[use as bounding box] (70.0,608.0) rectangle (548.0,766.0);
\end{tikzpicture}}\par



\par









\QHeading{第2、4题　答案不唯一的范围}

\Answer{梯形的高可以从一条底边上的不同点向另一底边所在直线作出，但必须与两底垂直。长方形、正方形、平行四边形、等腰梯形的自由作图可改变大小和位置；应核对各自的边、角条件，而不能只按本册示例的外观判断。}

\Answer{作高复核：高用虚线，旁边写“高”，垂足处画直角符号；以题目指定的底为准，必要时用虚线延长底边。}
